% Umlaute in der Datei erlauben, auf deutsch umstellen
\usepackage[utf8]{inputenc}
%\usepackage[ngerman]{babel}
\usepackage[nottoc,numbib]{tocbibind}
%\usepackage{mathptmx} % Hier steckt Times drin
%\usepackage{fontspec}
%\setmainfont{Times New Roman}
\usepackage{xcolor}


% Mathesymbole und Ähnliches
\usepackage{amsmath}
\usepackage{mathpartir}
\usepackage{mathtools}
\usepackage{amssymb}
\usepackage{amsthm}
\usepackage{mdframed}
%\usepackage{microtype}
\usepackage{stmaryrd}
\newcommand{\blitz}{\textcolor{red}{\lightning}}
\usepackage{float}
\renewcommand{\phi}{\varphi}


% Für Code
\usepackage{listings}
\lstset{frame=l,
  language=Java,
  basicstyle={\small\ttfamily},
  numberstyle=\tiny\color{blue},
  keywordstyle=\color{purple},
  commentstyle=\small\color{gray},
  stringstyle=\color{orange},
  breaklines=true,
  tabsize=2
}

% Pseudocode
\usepackage[vlined, linesnumbered, ruled]{algorithm2e}
\SetKwComment{Comment}{/* }{*/}
%\SetKwInput{KwData}{Input}
%\SetKwInput{KwResult}{Output}
\IncMargin{1.5em}

% Abbildungen
\usepackage{tikz}
\usetikzlibrary{arrows,calc,shapes,positioning}
\usepackage{wrapfig}
\usepackage[labelfont=bf]{caption}

% quotation
\definecolor{citeblue2}{RGB}{44,107,165}
\definecolor{citeblue}{rgb}{0,0.25,0.5}
\RequirePackage{csquotes}
\usepackage[backend=biber, style=numeric]{biblatex}
%\usepackage[citecolor=black, urlcolor=black, linkcolor=black]{hyperref}
\usepackage[citecolor=citeblue, urlcolor=citeblue, linkcolor=citeblue]{hyperref}
\hypersetup{colorlinks=true}
\usepackage[english, nameinlink, noabbrev]{cleveref}

% Custom Makros
\newcommand{\bangle}[1]{\langle #1 \rangle}
\newcommand{\semsof}[1]{\ensuremath{\llbracket #1 \rrbracket}}
\newcommand{\impl}{\Rightarrow}
\newcommand{\reason}[1]{\qquad \mid #1}
\newcommand{\edge}{\ensuremath{(e,e')}}
\newcommand{\edgein}{\ensuremath{(e,e') \in }}
\newcommand{\defof}{\text{Def. of }~}
\newcommand{\inv}{\ensuremath{{-1}}}
\newcommand{\prf}{\mathcal{P}}
\newcommand{\catprf}{\ensuremath{\mathcal{P}^{CAT}}}

% Reelle, Natürliche, Ganze, Rationale Zahlen
\newcommand{\R}{\ensuremath{\mathbb{R}}}
\newcommand{\N}{\ensuremath{\mathbb{N}}}
\newcommand{\Z}{\ensuremath{\mathbb{Z}}}
\newcommand{\Q}{\ensuremath{\mathbb{Q}}}
\newcommand{\E}{\ensuremath{\mathbb{E}}}
\newcommand{\B}{\ensuremath{\mathfrak{B}}}
\newcommand{\C}{\ensuremath{\mathfrak{C}}}


% Kürzel für Basisrelationen
\definecolor{rfgreen}{rgb}{0.2,0.55,0.2}
\newcommand{\0}{\textsf{\upshape 0}} % \upshape ensures that the relations remain their look in lemmata and inferrules
\newcommand{\po}{\textsf{\upshape po}}
\newcommand{\co}{\textsf{\textcolor{orange}{\upshape co}}}
\newcommand{\rf}{\textsf{\textcolor{rfgreen}{\upshape rf}}}
\newcommand{\ext}{\textsf{\upshape ext}}
\newcommand{\rmw}{\textsf{\upshape rmw}}
\newcommand{\hb}{\textsf{\textcolor{purple}{\upshape hb}}}
\newcommand{\loc}{\textsf{\upshape loc}}
\newcommand{\fr}{\textsf{\textcolor{violet}{\upshape fr}}}
\newcommand{\id}{\textsf{\upshape id}}
\newcommand{\true}{\textit{true}}
\newcommand{\false}{\textit{false}}
\newcommand{\rel}{\bangle{rel}}

%\newcommand{\rel}[1]{\textsf{#1}}
\newcommand{\relI}{\textsf{\upshape I}}
\newcommand{\relW}{\textsf{\upshape [W]}}
\newcommand{\relR}{\textsf{\upshape [R]}}
\newcommand{\relr}{\textsf{\upshape r}}
\newcommand{\Relr}{\relr}
\newcommand{\rels}{\textsf{\upshape s}}
\newcommand{\Rels}{\rels}
\newcommand{\relt}{\textsf{\upshape t}}


% Definitionsboxen
\usepackage[most]{tcolorbox}
%\definecolor{darkgray}{RGB}{70,90,70}
\definecolor{examblue}{rgb}{0,0.2,0.55}
\newenvironment{defi}[1]{%
\begin{tcolorbox}[blanker, before skip=5pt, after skip = 5pt, colback=darkgray!5!white, colframe=darkgray!75!black, title={{\sffamily} \textbf{#1}}]
}{
\end{tcolorbox}}

%Regelboxen
\tcbset{
  mybox/.style={
    colback=white,
    colframe=gray!70,
    colbacktitle=white,
    boxrule=2.5pt,
    coltitle=gray,
    arc=0pt,
    %before skip=0pt,
    %after skip=0pt,
    enhanced,
    left=6pt, right=6pt, top=6pt, bottom=6pt,
    attach boxed title to top left={yshift=-3mm, xshift=5mm},
    boxed title style={
    size=title,
    colback=white,
    colframe=white,
    boxrule=0.1pt
  }
  }
}
\newenvironment{rules}[1]{%
\begin{tcolorbox}[mybox, title={ \textbf{#1}}]
\begin{mathpar}
}{
\end{mathpar}
\end{tcolorbox}}

% Formelboxen
\newenvironment{formel}[1]{%
\begin{tcolorbox}[colback=white, colframe=violet!85!black, title={\sffamily \textbf{#1}}]
}{
\end{tcolorbox}}


%\theoremstyle{definition}
\newtheorem{theorem}{Theorem}%[section]

\newtheorem{proposition}[theorem]{Proposition}%[section]
\crefname{proposition}{Proposition}{Propositions}

\newtheorem{property}[theorem]{Eigenschaft}%[section]
\crefname{property}{Eigenschaft}{Eigenschaften}

\newtheorem{lemma}[theorem]{Lemma}%[section]
\crefname{lemma}{Lemma}{Lemmata}

%\newtheorem{definition}{Definition}
\mdfdefinestyle{theoremstyle}{%
linecolor=red,middlelinewidth=2pt,%
frametitlerule=true,%
frametitlebackgroundcolor=darkgray!60,
frametitlerulewidth=1pt,
innertopmargin=\topskip,
}
\mdfsetup{skipabove=\topskip , skipbelow=\topskip}
\mdfdefinestyle{defistyle}{linecolor=darkgray, linewidth=2pt, roundcorner=15pt, frametitlebackgroundcolor=gray!40, frametitlefont=\bfseries, theoremtitlefont=\bfseries, innertopmargin=5pt}
\mdtheorem[style=defistyle, nobreak=true]{definition}[theorem]{Definition}%[section]
\crefname{definition}{definition}{definitions}
\let\olddefinition\definition
\renewcommand{\definition}{\crefalias{theorem}{definition}\olddefinition}



% Beispiele
\newcounter{excount}
%\theoremstyle{definition}
\newtheorem{exi}[excount]{Beispiel}
\crefname{example}{Beispiel}{Beispiele}
\newenvironment{example}[2]{%
\begin{tcolorbox}[colback=white, colframe=examblue!85!black, title={\begin{exi}\label{#2} \textbf{#1}\end{exi}}]
}{
\end{tcolorbox}}




% Für Vektoren
\newcommand{\myvec}[1]{\ensuremath{\begin{pmatrix}#1\end{pmatrix}}}
\def\spvec#1{\left(\vcenter{\halign{\hfil$##$\hfil\cr \spvecA#1;;}}\right)}
\def\spvecA#1;{\if;#1;\else #1\cr \expandafter \spvecA \fi}

\usepackage{lipsum}

% PDFs includen
\usepackage{pdfpages}

\usetikzlibrary{positioning}
\tikzstyle{white} = [circle, minimum width=0.4cm, text centered, draw=black, fill=white]
\tikzstyle{blank} = [minimum width=0.6cm, text centered]
\tikzstyle{gray} = [circle, minimum width=0.6cm, text centered, text=white, draw=black, fill=darkgray]
\tikzstyle{blue} = [circle, minimum width=0.6cm, text centered, text=black, draw=black, fill=cyan]
\tikzstyle{purple} = [circle, minimum width=0.6cm, text centered, text=white, draw=black, fill=purple]
\tikzstyle{edge} = [thick,>=stealth]
\tikzstyle{arrow} = [thick,-stealth]
